\documentclass[a4paper,11pt]{ctexart}

\usepackage[top=1.8cm,bottom=1.8cm,left=2.1cm,right=2.1cm]{geometry}
\usepackage{amsmath,amssymb}
\usepackage{enumitem}
\usepackage{fancyhdr}
\usepackage{lastpage}

\pagestyle{fancy}
\fancyhf{}
\renewcommand{\headrulewidth}{0pt}
\fancyfoot[C]{\small 第~\thepage~页\quad 共~\pageref{LastPage}~页}

\linespread{1.22}
\setlength{\parindent}{0pt}
\setlength{\parskip}{0.18em}

\newcommand{\tr}{\operatorname{tr}}
\newcommand{\rk}{\mathrm{r}}
\newcommand{\T}{\mathrm{T}}

\begin{document}

\begin{center}
  {\Large\bfseries 《线性代数 2020 春 B》参考答案}
\end{center}

\vspace{0.6em}

\textbf{一、判断题}

1.\ 错，\quad
2.\ 对，\quad
3.\ 对，\quad
4.\ 错，\quad
5.\ 对，\quad
6.\ 对，\quad
7.\ 对，\quad
8.\ 对

\vspace{0.8em}

\textbf{二、填空题}

1.\ $2^{3n-1}$，\quad
2.\ $\alpha_5 = -\dfrac{9}{2}\alpha_1 + \dfrac{5}{2}\alpha_2 + \dfrac{1}{2}\alpha_4$，\quad
3.\ $48$，\quad
4.\ $\lambda \neq 1$，\quad
5.\ $(1, \infty)$

\vspace{0.5em}

\textbf{三、解答题}

解：按第一列展开
\[
  D_n = x D_{n-1} + (-1)^{n+1} a_n
  \begin{vmatrix}
    -1 & & & & \\
    x & -1 & & & \\
    & x & -1 & & \\
    & & \ddots & \ddots & \\
    & & & x & -1
  \end{vmatrix}
  = x D_{n-1} + a_n.
\]
所以递推得到
\[
  D_n = x^n + a_1 x^{n-1} + \cdots + a_{n-1} x + a_n.
\]

\vspace{0.5em}

\textbf{四、解答题}

方程有唯一解，当且仅当系数行列式非零，即
\[
  \begin{vmatrix}
    1 & -a & -2 \\
    1 & -1 & a \\
    5 & -5 & -4
  \end{vmatrix} \neq 0,
\]
解之得 $a \neq -\dfrac{4}{5}$ 且 $a \neq 1$.

当 $a = -\dfrac{4}{5}$ 时，系数矩阵与增广矩阵的秩不一致，所以无解；

当 $a = 1$ 时，无穷多解，其通解为 $c(1, 1, 0)^{\T} + (1, 0, 1)^{\T}$（其中 $c \in \mathbb{R}$）.

\vspace{0.5em}

\textbf{五、解答题} \par\nopagebreak

解：$f$ 所对应的矩阵为 $A = \begin{pmatrix} 1 & 2 & 2 \\ 2 & 1 & 2 \\ 2 & 2 & a \end{pmatrix}$. 由题意知 $A$ 的特征值为 $-1, -1, 5$. 于是根据迹的性质得到：$1 + 1 + a = -1 - 1 + 5$，所以 $a = 1$.

求特征值 $\lambda = -1$ 所对应的特征向量，即求解方程组 $(-I - A)X = 0$，解之得基础解系为
\[
  (-1, 1, 0)^{\T},\quad (-1, 0, 1)^{\T}.
\]
利用 Gram-Schmidt 正交化办法对其正交化和单位化得到：
\[
  \left(-\frac{1}{\sqrt{2}}, \frac{1}{\sqrt{2}}, 0\right)^{\T},\quad \left(-\frac{1}{\sqrt{6}}, -\frac{1}{\sqrt{6}}, \frac{2}{\sqrt{6}}\right)^{\T}.
\]
求特征值 $\lambda = 5$ 所对应的特征向量，即求解方程组 $(5I - A)X = 0$，解之得基础解系为 $(1, 1, 1)^{\T}$. 对其单位化得到：$\left(\dfrac{1}{\sqrt{3}}, \dfrac{1}{\sqrt{3}}, \dfrac{1}{\sqrt{3}}\right)^{\T}$.

所以将此二次型对应的标准形所采用的正交变换矩阵为
\[
  \begin{pmatrix}
    -\dfrac{1}{\sqrt{2}} & -\dfrac{1}{\sqrt{6}} & \dfrac{1}{\sqrt{3}} \\[0.8em]
    \dfrac{1}{\sqrt{2}} & -\dfrac{1}{\sqrt{6}} & \dfrac{1}{\sqrt{3}} \\[0.8em]
    0 & \dfrac{2}{\sqrt{6}} & \dfrac{1}{\sqrt{3}}
  \end{pmatrix}.
\]

\vspace{0.5em}

\textbf{六、}

解：$A$ 的特征值为 $3, 0, 0$，其中 $3$ 为 $1$ 重特征值，$0$ 为 $2$ 重特征值. 由于 $A$ 可相似对角化，所以 $0$ 有两个线性无关的特征向量（即 $(0I - A)X = 0$ 的解空间的维数为 $2$），从而 $\rk(A) = 1$，所以 $a = 0$.

\vspace{0.5em}

\textbf{七、}

解：根据 $f$ 的定义，得到：
\begin{align*}
  I + f(A) &= I + (I + A)^{-1}(I - A) = (I + A)^{-1}(I + A) + (I + A)^{-1}(I - A) \\
  &= (I + A)^{-1}(I - A + I + A) = 2(I + A)^{-1}.
\end{align*}
于是
\[
  \bigl(I + f(A)\bigr) \cdot \frac{I + A}{2} = I,
\]
所以 $I + f(A)$ 可逆，且其逆为 $\dfrac{I + A}{2}$.

\vspace{0.5em}

\textbf{八、}

证明：

必要性：考虑 $A$ 的行标准型. 由于 $A$ 的秩为 $n$，所以 $A$ 的行标准型中有 $n$ 列，分别是 $m$ 维的基本单位向量. 同时，又由于 $A$ 只有 $n$ 个列，所以 $A$ 的行标准型为 $\begin{pmatrix} I \\ O \end{pmatrix}$，其中 $I$ 是 $n$ 阶的单位矩阵，$O$ 是 $(m - n) \times n$ 零矩阵. 所以存在 $m$ 阶可逆矩阵 $G$ 使得 $GA = \begin{pmatrix} I \\ O \end{pmatrix}$. 最后，令 $B = (I, O^{\T})G$（其中 $O^{\T}$ 为 $n \times (m - n)$ 零矩阵），则
\[
  BA = (I, O^{\T})GA = (I, O^{\T})\begin{pmatrix} I \\ O \end{pmatrix} = I.
\]

充分性：若存在 $n \times m$ 矩阵 $B$ 使得 $BA = I$，则 $n = \rk(I) = \rk(BA) \le \rk(A) \le n$，故 $\rk(A) = n$.

\end{document}
